\begin{figure}[H]
\centering
\begingroup
\newcommand{\rootneighborhood}[2]{%
  \begin{tikzpicture}[x=6mm,y=6mm,baseline=-2.5pt]
    \coordinate (a) at (0,0);
    \coordinate (b) at (1,1);
    \coordinate (c) at (1,-1);
    \coordinate (d) at (.66,0);
    \coordinate (v) at (2.05,.20);
    \draw[line width=.55pt] (a)--(b)--(c)--cycle (a)--(d)--(b) (d)--(c);
    \foreach \endpoint in {#1} {\draw[densely dashed,line width=.45pt] (v)--(\endpoint);}
    \foreach \root in {a,b,c,d} {\fill (\root) circle[radius=1.15pt];}
    \draw[fill=white,line width=.55pt] (v) circle[radius=1.6pt];
  \end{tikzpicture}\,$#2$}
\setlength{\tabcolsep}{7pt}
\begin{tabular}{@{}cccc@{}}
\toprule
$q_v=1$ & $q_v=2$ & $q_v=3$ & $g$\\
\midrule
\rootneighborhood{b}{26} & \rootneighborhood{b,c}{27} & \rootneighborhood{b,c,d}{0} & $11$\\
\addlinespace[6pt]
\rootneighborhood{b}{25} & \rootneighborhood{b,c}{26} & \rootneighborhood{b,c,d}{1} & $7$\\
\addlinespace[6pt]
\rootneighborhood{b}{24} & \rootneighborhood{b,c}{25} & \rootneighborhood{b,c,d}{2} & $3$\\
\bottomrule
\end{tabular}
\endgroup
\caption{The three possible neighborhood counts at a marked $K_4$ for $D=23$.
The open vertex lies outside $A$. Dashed lines show its neighbors in $A$,
up to permutations of the four vertices.}
\label{fig:finite-k4-neighborhoods}
\end{figure}
